\documentclass[11pt]{article}

\usepackage{amsmath,amssymb}
\usepackage{xurl}
\usepackage[hidelinks]{hyperref}

% Shared commands for notes in docs/paper-gaps/.
% Notation follows the LDT paper (references/ldt-paper/).

\newcommand{\Lean}{\textsc{Lean}}
\newcommand{\leanid}[1]{\textsf{#1}}
\newcommand{\ghissue}[1]{\href{https://github.com/LionSR/MIPStarRE/issues/#1}{GitHub issue \##1}}
\newcommand{\ghpr}[1]{\href{https://github.com/LionSR/MIPStarRE/pull/#1}{GitHub PR \##1}}

% --- Paper-compatible notation ---
\newcommand{\F}{\mathbb{F}}
\newcommand{\N}{\mathbb{N}}
\newcommand{\C}{\mathbb{C}}
\newcommand{\tr}{\operatorname{tr}}
\newcommand{\ot}{\otimes}
\newcommand{\E}{\mathop{\bf E\/}}
\newcommand{\bx}{\mathbf{x}}
\newcommand{\by}{\mathbf{y}}
\newcommand{\bu}{\mathbf{u}}
\newcommand{\bv}{\mathbf{v}}

% Approximation relation (paper uses \approx_\eps and \simeq_\zeta)
\newcommand{\approxeps}{\approx_{\varepsilon}}
\newcommand{\simgeq}{\simeq_{\zeta}}

% Polynomial spaces (paper uses \polyfunc, \polysub, \polymeas)
\newcommand{\polyfunc}[3]{\operatorname{Poly}(#1,#2,#3)}
\newcommand{\polysub}[3]{\operatorname{PolySub}(#1,#2,#3)}
\newcommand{\polymeas}[3]{\operatorname{PolyMeas}(#1,#2,#3)}


\title{Issue 904: The Missing Completion Term in the Definition of
\texorpdfstring{$\zeta_2$}{zeta2}}
\author{}
\date{2026-04-29}

\newcommand{\dd}{\delta}
\newcommand{\zz}{\zeta}

\begin{document}
\maketitle

\section{The estimate for the completion step}

This note concerns the scalar in the completion estimate in
\cite[Section~\emph{The main induction step}]{Ji2020LowIndividualDegree}.
\footnote{The issue record is
\href{https://github.com/LionSR/MIPStarRE/issues/904}{GitHub issue \#904}; the
corresponding repair is
\href{https://github.com/LionSR/MIPStarRE/pull/909}{GitHub pull request \#909}.
In the TeX source consulted here, the estimate occurs in
\path{references/ldt-paper/inductive_step.tex}, lines 135--149.}
We write \(X\approx_\eta Y\) for the state-dependent squared-distance relation,
with respect to the shared state \(\psi\) fixed in the surrounding induction
argument, and with error parameter \(\eta\).  The symbol \(I\) denotes the
identity operator on the tensor factor which is not being measured.

At this point in the induction, \(G^{\mathrm A}\) and \(G^{\mathrm B}\) are the
polynomial-outcome measurements before projectivization.  Their
self-consistency error, obtained in the preceding part of the induction, is
denoted by \(\zeta_1\).  Applying the orthogonalization lemma for measurements to
this self-consistency estimate gives, for each polynomial outcome \(g\),
\begin{align*}
  G^{\mathrm A}_g \otimes I
    &\approx_{100\zeta_1^{1/4}}
  P^{\mathrm A}_g \otimes I,
    \tag{\path{lem:orthonormalization-main-lemma}}\\
  I \otimes G^{\mathrm B}_g
    &\approx_{100\zeta_1^{1/4}}
  I \otimes P^{\mathrm B}_g.
    \tag{\path{eq:G-self-consistency}}
\end{align*}
Here \(P^{\mathrm A}\) and \(P^{\mathrm B}\) are projective submeasurements.  The
next step is to complete
these submeasurements to projective measurements \(Q^{\mathrm A}\) and
\(Q^{\mathrm B}\).  The value assigned there to the new error parameter is
\(\zeta_2=200\zeta_1^{1/4}+40\zeta_1^{1/8}\), and the asserted estimate is
\begin{align*}
  G^{\mathrm A}_g \otimes I
    &\approx_{\zeta_2}
  Q^{\mathrm A}_g \otimes I,\\
  I \otimes G^{\mathrm B}_g
    &\approx_{\zeta_2}
  I \otimes Q^{\mathrm B}_g.
    \tag{\path{eq:G-with-Q-A}}
\end{align*}

\section{The completion calculation}

The needed completion result is the completion proposition in
\cite[Section~\emph{Preliminaries}]{Ji2020LowIndividualDegree}.
\footnote{In the TeX source consulted here, this proposition occurs in
\path{references/ldt-paper/preliminaries.tex}, lines 1101--1111, and has source label
\path{prop:completing-to-measurement}.}
In the special case needed here, \(A=\{A_a\}\) is a measurement, \(B=\{B_a\}\)
is a submeasurement, and \(C=\{C_a\}\) is obtained from \(B\) by adding the
missing mass \(I-\sum_a B_a\) to one distinguished outcome.  The proposition
says that if
\[
  A_a\otimes I \approx_{\dd} B_a\otimes I
\]
and the measurement \(A\) has self-consistency error \(\zz\), then
\[
  A_a\otimes I
    \approx_{2\dd+4\sqrt{\dd}+2\zz}
  C_a\otimes I.
\]

For the projectivization step, \(A\) is the \(G\)-measurement, \(B\) is the
projective submeasurement \(P\), and \(C\) is the completed measurement \(Q\).
Thus \(\dd=100\zeta_1^{1/4}\), while the self-consistency parameter entering
the completion proposition is \(\zz=\zeta_1\).  Substitution gives
\begin{align*}
  2\dd+4\sqrt{\dd}+2\zz
    &=
  2(100\zeta_1^{1/4})
    +4\sqrt{100\zeta_1^{1/4}}
    +2\zeta_1\\
    &=
  200\zeta_1^{1/4}
    +40\zeta_1^{1/8}
    +2\zeta_1 .
\end{align*}
Thus the scalar \(\zeta_2\) used in the completion estimate omits the final term
\(2\zeta_1\).  This term is not a matter of notation: it is the contribution of
the self-consistency error in the completion proposition.

\section{Absorbing the missing term}

In the non-vacuous branch of the cascade, the standing small-error hypotheses
give \(0\leq\zeta_1\leq 1\).  In this range the omitted term is harmless at the
scale already being used.  Indeed \(x\leq x^{1/8}\) for \(0\leq x\leq 1\), and
hence \(2\zeta_1\leq 2\zeta_1^{1/8}\).  The literal completion error is
therefore bounded by \(200\zeta_1^{1/4}+42\zeta_1^{1/8}\).

Equivalently, the formal development keeps the exact composed error
\(2(100\zeta^{1/4})+4\sqrt{100\zeta^{1/4}}+2\zeta\) and proves, for
\(0\leq \zeta\leq 1\), the bound
\[
  2(100\zeta^{1/4})+4\sqrt{100\zeta^{1/4}}+2\zeta
  \leq 200\zeta^{1/4}+42\zeta^{1/8}.
\]
Applying this with \(\zeta=\zeta_1\) gives the estimate needed for the
completion step.\footnote{The corresponding Lean declarations are
\path{orthonormalizeAndCompleteError_le_absorbedZeta2},
\path{cascadeZeta2}, and
\path{orthonormalizeAndCompleteError_zeta1_le_zeta2}.  They are implemented in
\path{MIPStarRE/LDT/MakingMeasurementsProjective/ProjectivizationChain.lean} and
\path{MIPStarRE/LDT/Test/ErrorCascade.lean}.}

Thus the scalar denoted \(\zeta_2\) in the corrected cascade used in the formal
proof is
\(\zeta_2=200\zeta_1^{1/4}+42\zeta_1^{1/8}\), rather than the scalar with
coefficient \(40\) in the completion estimate.

\section{Relation with the blueprint and Lean}

The blueprint follows the corrected scalar
\cite[Chapter~\emph{The main induction step}]{MIPStarREBlueprint}.  Its cascade
definition uses
\(\zeta_2=200\zeta_1^{1/4}+42\zeta_1^{1/8}\), and explains that the coefficient
\(42\) absorbs the residual term \(2\zeta_1\).\footnote{In the blueprint source,
this is in \path{blueprint/src/chapter/ch10_induction.tex}, lines 433--451, with
source label \path{def:main-formal-error-cascade}.}
Later in the same chapter, the informal proof of the projectivization step again
identifies the literal completion error as
\[
  2(100\zeta_1^{1/4})+4\sqrt{100\zeta_1^{1/4}}+2\zeta_1.
\]
Thus the blueprint agrees with the checked scalar calculation, rather than with
the coefficient \(40\) in
\cite[Section~\emph{The main induction step}]{Ji2020LowIndividualDegree}.

The scalar correction is checked in the projectivization chain and in the final
scalar cascade.  These two parts of the development contain the literal
completion error, the absorption into the coefficient \(42\), and the downstream
cascade bounds.

The unfinished proof of the final theorem is not part of this scalar
calculation.\footnote{At the date of this note, that unrelated unfinished proof is in
\path{MIPStarRE/LDT/Test/MainTheorem.lean}.}  The calculation in
\href{https://github.com/LionSR/MIPStarRE/issues/904}{issue \#904} has been
isolated and accounted for in the checked scalar and
projectivization lemmas.

\section{Conclusion}

The verdict is that
\href{https://github.com/LionSR/MIPStarRE/issues/904}{issue \#904} records a
real arithmetic omission in an intermediate error parameter, but not an
obstruction arising from this scalar calculation to the final soundness estimate.

The formula for \(\zeta_2\) used in the completion estimate lacks the term
\(2\zeta_1\) arising from the completion proposition.  The corrected scalar
\(200\zeta_1^{1/4}+42\zeta_1^{1/8}\) absorbs this term in the regime
\(0\leq\zeta_1\leq 1\), and is still bounded by the error parameter appearing
in the final soundness theorem.  The discrepancy is therefore a quantitative
error in an intermediate scalar, not by itself an obstruction to the final
theorem.

\bibliographystyle{alpha}
\bibliography{docs/paper-gaps/references}

\end{document}
